% Certified Fuzzy Control by Least Action
% Compile: pdflatex least_action_fis.tex  (twice, for references)
% For NAFIPS submission, swap \documentclass for \documentclass[conference]{IEEEtran}
% and delete the \geometry line.
\documentclass[11pt]{article}

\usepackage[margin=1in]{geometry}
\usepackage{amsmath,amssymb,amsthm}
\usepackage{booktabs}
\usepackage{array}
% microtype's font expansion needs scalable fonts; omitted so the source
% compiles on a minimal TeX installation.
\usepackage[hidelinks]{hyperref}

\theoremstyle{plain}
\newtheorem{theorem}{Theorem}
\newtheorem{proposition}[theorem]{Proposition}
\newtheorem{corollary}[theorem]{Corollary}
\theoremstyle{definition}
\newtheorem{remark}[theorem]{Remark}

\newcommand{\R}{\mathbb{R}}
\newcommand{\Hone}{H^1}
\newcommand{\inner}[2]{\langle #1,#2\rangle_{\Hone}}
\newcommand{\Vdot}{\dot{V}}

\title{\textbf{Certified Fuzzy Control by Least Action}\\[2mm]
\large An $H^1$ Variational Formulation for TSK Systems, with\\
Sum-of-Squares Stability Certificates and a Calibrated\\
Performance/Certificate Trade}

\author{
\begin{tabular}{c@{\hspace{1.2cm}}c}
Scott Phillips & Kelly Cohen \\
\end{tabular}\\[2mm]
\small Department of Aerospace Engineering and Engineering Mechanics\\
\small University of Cincinnati\\[1mm]
\small \texttt{polygonguru@gmail.com}}

\date{\today}

\begin{document}
\maketitle

\begin{abstract}
We develop a variational formulation of Takagi--Sugeno--Kang (TSK) fuzzy
inference in which model fitting is posed as minimization of an $H^1$
(Sobolev) action functional, and we carry the formulation through to
closed-loop control with machine-checked stability certificates. Four results
are analytic. First, with antecedents fixed the consequent block is the
\emph{global} minimizer of the action in closed form, so the only non-convexity
in the problem is the antecedent placement; variable projection then yields an
exact reduced gradient via the envelope theorem. Second, the natural design
question---when does sequential (greedy) rule identification coincide with the
joint fit---has the exact answer ``iff the $H^1$ Gram matrix is
block-diagonal,'' and for two equal-width Gaussians the slope weight that
achieves $H^1$ orthogonality has the closed form
$\lambda^{*}=3b^{4}/(8c^{2})$, equivalently $\ell^{*}=\sqrt{6}\,w$ where $w$ is
the crossover width. A factor-of-four obstruction shows no single $\lambda$ can
orthogonalize three or more uniformly spaced rules. Third, on the control side
the closed-loop suboptimality gap is \emph{exactly} an integral of a Bregman
divergence for any convex control cost, not merely the quadratic case. Fourth,
switching to $\pi$-shaped (rational) membership functions makes the closed loop
a rational vector field, and a sum-of-squares S-procedure certificate for the
region of attraction is then verified in exact rational arithmetic, with no
floating-point step in the chain.

Empirically, on a two-mass-spring benchmark scored jointly on settling time,
peak force and energy, the least-action fit converges in two rules ($10$
parameters, score $0.867$ against an LQR baseline of $1.000$) and then degrades
through distribution shift rather than through any deficiency of the fit.
Direct policy optimization warm-started from the least-action fit reaches
$0.655$, improving all three objectives simultaneously, while a cold start
fails outright---the variational fit is the initializer that makes the
optimizer reachable. Finally we make the stability certificate an
\emph{objective} rather than an afterthought. A sampled necessary condition for
the sum-of-squares program, evaluated by ray search, estimates the certified
region $2500\times$ faster than the semidefinite program with a residual bias
of a single constant $\kappa=0.962$; specifying a certified radius then yields
a controller the solver certifies at or above it. The best controller so
obtained carries a certified ball of $0.679$ at score $0.847$, dominating the
imitation fit on both axes.
\end{abstract}

\tableofcontents

\section{Introduction}

This paper develops an idea originating in a handwritten derivation by the
second author: that fitting a fuzzy inference system should be posed as a
principle of least action rather than as least squares, and that the resulting
stationarity conditions should be strong enough to certify optimality rather
than merely to terminate an optimizer. The
motivation is that a TSK system is a smooth parametric model, and smooth
parametric models admit variational treatment; what is unusual is insisting
that every step of the treatment produce a \emph{statement that can be
checked}, rather than a number produced by a solver.

Two questions organize the work.

\begin{enumerate}
\item \textbf{Approximation.} Given a target field $y_d$ on an input universe
$X$, what is the right functional to minimize over TSK parameters, and what can
be proved about its minimizers? In particular: under what conditions does
sequential rule identification (fit rule 1, fit rule 2 to the residual, and so
on) agree with the joint fit? This was the question that prompted the work.

\item \textbf{Control.} If the same machinery is turned on a control problem,
what is provable about the resulting controller---not measured, proved? We ask
for suboptimality certificates and for stability certificates, and we ask what
happens when the two are placed in competition.
\end{enumerate}

\paragraph{Contributions.} The specific claims are:
\begin{itemize}
\item An $H^1$ action functional for TSK approximation whose consequent block
has a closed-form global minimizer (Theorem~\ref{thm:varpro}), an exact
analytic reduced gradient (Theorem~\ref{thm:grad}), and second-order
conditions stated correctly on the critical cone (\S\ref{sec:second-order}).
\item An exact characterization of when greedy identification equals the joint
fit (Proposition~\ref{prop:greedy}), a closed form for the $H^1$-orthogonalizing
slope weight (Theorem~\ref{thm:lambda}), and a structural obstruction showing it
cannot be achieved for $N\ge3$ uniformly spaced rules
(Corollary~\ref{cor:three}).
\item The classifier-side results that make defuzzification differentiable:
Mamdani aggregation equals summation iff output supports are disjoint
(Theorem~\ref{thm:disjoint}), centroid defuzzification of such a system is an
order-$0$ TSK model (Theorem~\ref{thm:centroid}), and an explicit error bound
for $\beta$-annealed mean-of-maxima (Theorem~\ref{thm:anneal}).
\item An exact suboptimality certificate for arbitrary convex control cost,
expressed as an integral of a Bregman divergence (Theorem~\ref{thm:c2prime});
the familiar quadratic identity is a special case.
\item A sum-of-squares region-of-attraction certificate for a $\pi$-membership
TSK closed loop, verified in exact rational arithmetic
(\S\ref{sec:sos}), together with the practical finding that the binding
obstruction was solver accuracy and not the relaxation.
\item A certificate-aware policy optimization in which the certified radius
enters the objective (\S\ref{sec:aware}), and a calibration procedure that turns
the certified radius from a dial into a specification (\S\ref{sec:calibrate}).
\end{itemize}

We are explicit throughout about what is proved, what is certified numerically,
and what is merely measured; \S\ref{sec:inventory} tabulates this.

\section{Model and action functional}
\label{sec:model}

\subsection{The model}

On a scalar input universe $X=[x_{lo},x_{hi}]$, an order-$p$ TSK system with
$N$ rules is
\begin{equation}
y_c(x)=\sum_{i=1}^{N}\varphi_i(x)f_i(x),\qquad
\varphi_i=\frac{\mu_i}{\sum_{j}\mu_j},\qquad
f_i(x)=\sum_{k=0}^{p}B_{ik}x^{k},
\label{eq:tsk}
\end{equation}
where $\mu_i$ is the membership function of rule $i$ with centre $a_i$ and
width $b_i$. The normalized firing strengths form a partition of unity,
$\sum_i\varphi_i\equiv1$. Writing $\psi_{ik}=\varphi_i x^{k}$ for the
\emph{rule regressors} and $\theta$ for the flattened coefficients $B$, the
model is linear in $\theta$:
\begin{equation}
y_c(x)=\theta^{\top}\psi(x).
\label{eq:linear}
\end{equation}
This is the structural fact the entire paper exploits. The parameters split
into a block in which the model is linear ($\theta$) and a block in which it is
not ($a,b$).

\begin{remark}
In implementation the monomial $x^k$ is replaced by a rescaled monomial
$t=(x-\mathrm{mid})/\mathrm{half}$. This spans the identical function space, so
the model is unchanged, but it prevents the Gram condition number from reaching
$10^{7}$ on a universe such as $[-15,15]$.
\end{remark}

\subsection{The action}

Let $e=y_d-y_c$ be the error field. We take as action the squared $H^1$ norm of
the error,
\begin{equation}
S[y_c]=\int_X\!\Big(e^2+\lambda\,(e')^2\Big)\,dx
=\|e\|_{H^1}^{2},\qquad \lambda=\ell^{2}\ge0,
\label{eq:action}
\end{equation}
where $\ell$ has units of $x$ and plays the role of a correlation length. The
$L^2$ term penalizes pointwise error; the $\lambda$ term penalizes error in
slope, which is what makes the formulation a genuinely variational one rather
than a reweighted regression.

\paragraph{Euler--Lagrange.} Treating $y_c$ as a free field, the first variation
of \eqref{eq:action} vanishes when
\begin{equation}
\lambda e''-e=0,
\label{eq:el}
\end{equation}
the screened-Poisson (modified Helmholtz) equation, with natural boundary
conditions $e'=0$ at the endpoints. Its only solution decaying in both
directions is $e\equiv0$: in the unconstrained function space, least action
recovers the target exactly. The content of the method is therefore entirely in
what happens when $y_c$ is restricted to the TSK manifold \eqref{eq:tsk}.

\paragraph{Stationarity in the model class.} Restricting to \eqref{eq:linear},
$\delta S=0$ in the $\theta$ directions is
\begin{equation}
\inner{e}{\psi_m}=0\quad\text{for every regressor }\psi_m,
\label{eq:galerkin}
\end{equation}
i.e.\ \emph{Galerkin orthogonality}: the error is $H^1$-orthogonal to the model
subspace. Equation~\eqref{eq:galerkin} is the weak form of \eqref{eq:el} and is
the correct stationarity condition for the finite-dimensional problem. It is
what we check numerically, and it is what distinguishes this formulation from a
generic nonlinear least-squares fit.

\begin{remark}[Convexity]
For $\lambda\ge0$ the map $y_c\mapsto S$ is a strictly convex quadratic on
function space, so the variational problem has no conjugate points and the
Jacobi condition is vacuous. All non-convexity in the parametric problem enters
through the map $(a,b)\mapsto\psi$, not through the action.
\end{remark}

\section{The consequent block is globally solvable}
\label{sec:consequents}

\subsection{Closed form}

Define the $H^1$ Gram matrix and right-hand side
\begin{equation}
G_{mn}=\inner{\psi_m}{\psi_n},\qquad r_m=\inner{y_d}{\psi_m}.
\end{equation}

\begin{theorem}[Variable projection]
\label{thm:varpro}
For fixed antecedents $(a,b)$, the action \eqref{eq:action} is a convex
quadratic in $\theta$ with Hessian $2G\succeq0$, and every minimizer satisfies
the normal equations $G\theta=r$. If the regressors are linearly independent in
$H^1$ then $G\succ0$ and $\theta^{*}=G^{-1}r$ is the unique global minimizer.
The reduced action is
\begin{equation}
\widehat S(a,b)=\|y_d\|_{H^1}^{2}-r^{\top}G^{-1}r.
\label{eq:reduced}
\end{equation}
\end{theorem}

\begin{proof}
Substituting \eqref{eq:linear} into \eqref{eq:action},
$S=\|y_d\|^2_{H^1}-2\theta^\top r+\theta^\top G\theta$. The Hessian $2G$ is a
Gram matrix of an inner product and hence positive semidefinite, so $S$ is
convex and stationarity is sufficient. Stationarity is $G\theta=r$, which is
\eqref{eq:galerkin} written out. Substituting $\theta^{*}$ gives
\eqref{eq:reduced}.
\end{proof}

Theorem~\ref{thm:varpro} is the strongest guarantee in the paper and it is
free: no iteration, no initialization, no local minima. It also localizes the
difficulty precisely---the problem is non-convex \emph{only} in $(a,b)$.

\begin{remark}[Regularization]
Numerically we solve $(G+\varepsilon\,\mathrm{tr}(G)/m\,I)\theta=r$ with
$\varepsilon=10^{-9}$. The ridge is trace-relative so it scales with the
problem, and it is a ridge rather than a truncating pseudo-inverse
deliberately: a rank-flipping SVD makes the map $(a,b)\mapsto\theta$
discontinuous, which destroys the differentiability the next section depends on.
\end{remark}

\subsection{Exact reduced gradient}

Finite differences on \eqref{eq:reduced} are unreliable, because $\widehat S$ is
computed only to accuracy $\kappa(G)\,\varepsilon_{\text{mach}}$; in our
experiments SciPy's default step produced a gradient dominated by round-off, and
L-BFGS-B terminated after zero iterations reporting success. The envelope
theorem removes the issue.

\begin{theorem}[Analytic reduced gradient]
\label{thm:grad}
Let $\theta^{*}(a,b)$ be as in Theorem~\ref{thm:varpro}. Then for any
antecedent parameter $p\in\{a_i,b_i\}$,
\begin{equation}
\frac{\partial\widehat S}{\partial p}
=-2\Big\langle e,\ \frac{\partial y_c}{\partial p}\Big\rangle_{H^1},
\qquad
\frac{\partial y_c}{\partial p}=\theta^{*\top}\frac{\partial\psi}{\partial p},
\label{eq:grad}
\end{equation}
with no contribution from $\partial\theta^{*}/\partial p$.
\end{theorem}

\begin{proof}
$\widehat S(a,b)=S(\theta^{*}(a,b),a,b)$, so by the chain rule
$\partial_p\widehat S=\partial_\theta S\cdot\partial_p\theta^{*}
+\partial_p S$. At $\theta^{*}$ the first factor vanishes by
Theorem~\ref{thm:varpro}. The remaining partial is \eqref{eq:grad}.
\end{proof}

The derivative $\partial\psi/\partial p$ is available in closed form from the
membership-function derivatives $(\mu,\partial_x\mu,\partial_a\mu,\partial_b\mu,
\partial^2_{xa}\mu,\partial^2_{xb}\mu)$; the $\partial_x$ terms are required
because the $H^1$ inner product differentiates the regressors. Replacing finite
differences by \eqref{eq:grad} was the difference between an optimizer that
stalled immediately and one that converged.

\subsection{Second-order conditions, stated correctly}
\label{sec:second-order}

Three errors are easy to make here and all three occurred during this work.

\begin{enumerate}
\item \textbf{Test the reduced Hessian, not the full one.} The full-parameter
Hessian is never definite: consequent rescaling and rule relabeling guarantee
flat and symmetric directions. Only $\widehat S$ has a meaningful Hessian.

\item \textbf{Fitted solutions sit on identifiability constraints.} We impose a
minimum rule gap ($0.6\times$ pitch) and width bounds ($[0.25,1.5]\times$
pitch) so that rules remain distinguishable. At an active constraint the
correct second-order condition is definiteness on the \emph{critical cone}---the
null space of the active constraint normals---not on all of $\R^{2N}$; and the
correct first-order condition is that the \emph{projected} gradient vanishes,
the raw gradient being balanced by KKT multipliers.

\item \textbf{Curvature is informative only at a stationary point.} Galerkin
orthogonality \eqref{eq:galerkin} certifies the consequent block by
construction and therefore says nothing about $(a,b)$; the antecedent gradient
must be checked separately.
\end{enumerate}

Our certificate accordingly reports the KKT residual on the critical cone, the
projected-Hessian spectrum, and a falsification probe that walks along the
most-negative eigenvector with the first-order term subtracted, to test whether
the action actually decreases. Because the Hessian is obtained by finite
differences, the verdict is tolerance-based: we compare against an estimated
noise floor $\varepsilon_f/h^{2}$ rather than against exact zero, since judging
against zero reads numerical dust as a saddle. This is the one place where we
certify \emph{numerically} rather than analytically (guarantee G11 in
\S\ref{sec:inventory}).

\section{Rule orthogonality and sequential identification}
\label{sec:orthogonality}

The original design question was whether rules can be constrained so that
$\int f_i f_j\,dx=0$ for $i\ne j$, and whether rules can then be identified
sequentially from residuals. Both halves have exact answers, and they are not
the expected ones.

\subsection{Greedy equals joint iff the Gram is block-diagonal}

Partition the regressor index set by rule, and let $G$ be blocked accordingly.

\begin{proposition}
\label{prop:greedy}
Sequential fitting---fit rule $1$, then fit rule $2$ to the residual, and so
on---returns the joint minimizer of \eqref{eq:action} for \emph{every} target
$y_d$ if and only if $G$ is block-diagonal with respect to the rule partition.
\end{proposition}

\begin{proof}
If $G$ is block-diagonal the normal equations decouple into one system per rule
and greedy solves each exactly. Conversely, suppose $G_{IJ}\ne0$ for some
$I\ne J$. Choose $y_d$ in the span of block $J$'s regressors. Greedy assigns
block $I$ a nonzero coefficient (since $r_I=G_{IJ}\theta_J\ne0$), whereas the
joint solution assigns block $I$ zero; the two differ.
\end{proof}

Note the quantifier: block-diagonality is required for agreement on
\emph{all} targets. For a particular $y_d$ greedy may coincide with the joint
fit accidentally.

\begin{proposition}
\label{prop:notdisjoint}
For $\lambda>0$, block-diagonality of $G$ does \emph{not} require disjoint rule
supports.
\end{proposition}

\begin{proof}[Proof sketch]
The block entry is
$\int\varphi_i\varphi_j\,dx+\lambda\int\varphi_i'\varphi_j'\,dx$. For adjacent
normalized rules the first integral is strictly positive and the second is
strictly negative, because the partition of unity forces
$\varphi_i'=-\varphi_j'$ in the two-rule crossover. Cancellation is therefore
possible at a strictly positive $\lambda$ with overlapping supports.
\end{proof}

Proposition~\ref{prop:notdisjoint} corrects a claim we initially made in the
other direction, obtained by setting $\lambda=0$ in a statement whose hypothesis
required $\lambda>0$.

\subsection{The orthogonalizing slope weight in closed form}

Proposition~\ref{prop:notdisjoint} raises the constructive question: what value
of $\lambda$ achieves cancellation? For the two-rule case it is elementary.

\begin{theorem}[$\lambda^{*}$]
\label{thm:lambda}
For two equal-width Gaussian membership functions with centres $\pm c$ and
width $b$ on $\R$, the normalized weights collapse to a logistic,
$\varphi_1=\sigma(kx)$ with $k=4c/b^{2}$, and
\begin{equation}
\int_{\R}\varphi_0\varphi_1\,dx=\frac1k,\qquad
\int_{\R}\varphi_0'\varphi_1'\,dx=-\frac{k}{6},
\end{equation}
so that the unique $\lambda$ making the two rules $H^1$-orthogonal is
\begin{equation}
\boxed{\ \lambda^{*}=\frac{6}{k^{2}}=\frac{3b^{4}}{8c^{2}}\ }
\qquad\Longleftrightarrow\qquad
\ell^{*}=\sqrt{\lambda^{*}}=\sqrt{6}\,w,\quad w=\frac{1}{k}=\frac{b^{2}}{4c}.
\end{equation}
\end{theorem}

\begin{proof}
The log-ratio $\log(\mu_1/\mu_0)=4cx/b^{2}$ is linear in $x$, hence
$\varphi_1=\sigma(kx)$ with $k=4c/b^{2}$ and
$\varphi_0\varphi_1=\sigma(1-\sigma)=\sigma'(kx)$, $\varphi_i'=\pm k\sigma'$.
Substituting $u=kx$ and then $s=\sigma(u)$, $ds=s(1-s)du$:
$\int\varphi_0\varphi_1dx=k^{-1}\int_\R\sigma'(u)du=k^{-1}$ and
$\int\varphi_0'\varphi_1'dx=-k\int_\R\sigma'(u)^2du=-k\int_0^1 s(1-s)ds=-k/6$.
Setting $k^{-1}-\lambda k/6=0$ gives the result.
\end{proof}

\paragraph{Interpretation.} The target $y_d$ appears nowhere in
Theorem~\ref{thm:lambda}. Therefore $\lambda^{*}$ is \emph{not} a distinguished
correlation length of the data; it is the correlation length of the rule
crossover region, fixed entirely by partition geometry. The ratio
$\ell^{*}/w=\sqrt{6}\approx2.44949$ was confirmed numerically in every
configuration tested.

\begin{corollary}[Why $N\ge3$ fails]
\label{cor:three}
On a uniform partition of pitch $d$, adjacent pairs (separation $d$) and
next-nearest pairs (separation $2d$) demand orthogonalizing weights in the
ratio $(2d/d)^{2}=4$, since $\lambda^{*}\propto c^{-2}$. A single scalar
$\lambda$ therefore cannot orthogonalize three or more uniformly spaced rules.
\end{corollary}

The measured ratio was $4.0000$ in every case, and this explains why the
residual off-block coupling plateaus near $10^{-2}$ for $N\ge3$ rather than
approaching zero.

\paragraph{The price of $\lambda^{*}$.} Choosing $\lambda=\lambda^{*}$ trades
$L^2$ accuracy for decoupling. Table~\ref{tab:price} measures that trade
against the $L^2$-optimal ($\lambda=0$) fit of the same partition, target
$\tanh(x/3)$.

\begin{table}[t]
\centering
\caption{Cost of imposing $H^1$ orthogonality, relative to the $L^2$-optimal
fit of the same partition.}
\label{tab:price}
\begin{tabular}{cccccccc}
\toprule
$c$ & $b$ & $w$ & $\lambda^{*}$ & $\ell^{*}$ &
$L^2$ at $\lambda{=}0$ & $L^2$ at $\lambda^{*}$ & cost \\
\midrule
5 & 1 & 0.050 & 0.015 & 0.12 & 3.2540 & 3.2542 & 0.01\% \\
5 & 2 & 0.200 & 0.240 & 0.49 & 1.0992 & 1.1036 & 0.40\% \\
5 & 4 & 0.800 & 3.840 & 1.96 & 0.5191 & 0.5394 & 3.91\% \\
5 & 6 & 1.800 & 19.44 & 4.41 & 0.1799 & 0.1938 & 7.72\% \\
5 & 8 & 3.200 & 61.44 & 7.84 & 0.7769 & 0.8204 & 5.60\% \\
3 & 4 & 1.333 & 10.67 & 3.27 & 0.1101 & 0.1174 & 6.59\% \\
\bottomrule
\end{tabular}
\end{table}

The cost stays below $8\%$ across the sweep and below $2\%$ for sharp
crossovers. The closed form explains why: $\ell^{*}=\sqrt6\,b^{2}/(4c)$ is
sub-rule-scale for any sensible partition, and weighting slopes at a short
correlation length barely perturbs the $L^2$ solution.

\subsection{The practical alternative}

Since exact orthogonality is unavailable for $N\ge3$, the usable route is to
orthogonalize the \emph{regressors} rather than the rules: an $H^1$
Gram--Schmidt (orthogonal least squares) pass on $\psi$, selecting rules
greedily by residual reduction. This restores the greedy/joint agreement of
Proposition~\ref{prop:greedy} by construction, at the cost of rules that are no
longer individually interpretable in the original basis.

\section{Classifier case: making defuzzification differentiable}
\label{sec:classifier}

Without recurrence, a classifier ranking is a regression model with range
$[0,1]$, so the same action applies. The obstruction is that the standard
defuzzifications---De Baets' first/last/mean of maxima---are discontinuous, and
a discontinuous output has no action to minimize. Three results remove the
obstruction.

\begin{theorem}[Aggregation equals summation]
\label{thm:disjoint}
Let rules have output membership functions $\{\nu_i\}$ on $Y$. Then
$\bigvee_i\alpha_i\nu_i=\sum_i\alpha_i\nu_i$ pointwise for all firing strengths
$\alpha$ if and only if the supports of $\{\nu_i\}$ are pairwise disjoint.
\end{theorem}

\begin{proof}
If supports are disjoint, at most one term is nonzero at each $y$, so max and
sum agree. Conversely if $\nu_i(y)\nu_j(y)>0$ for some $y$ and $i\ne j$, take
$\alpha_i=\alpha_j=1$: the sum strictly exceeds the max.
\end{proof}

\begin{theorem}[Centroid is order-$0$ TSK]
\label{thm:centroid}
Under the hypotheses of Theorem~\ref{thm:disjoint} with product implication,
centroid defuzzification returns
$\hat y=\sum_i\varphi_i\,\bar y_i$ where $\bar y_i$ is the centroid of $\nu_i$
and $\varphi_i$ the normalized firing strength; i.e.\ exactly an order-$0$ TSK
model, whose output range $[0,1]$ is free rather than imposed.
\end{theorem}

\begin{theorem}[Annealed mean-of-maxima bound]
\label{thm:anneal}
Let $\hat y_\beta$ be the softmax-annealed ($\beta$-weighted) defuzzification
and $\hat y_\infty$ the exact mean-of-maxima. If $D$ is the diameter of the
output universe and $r<1$ is the ratio of the second-largest to the largest
firing strength, then
\begin{equation}
|\hat y_\beta-\hat y_\infty|\ \le\
\frac{(N-1)\,D\,r^{\beta}}{1+(N-1)r^{\beta}} .
\end{equation}
\end{theorem}

Theorem~\ref{thm:anneal} makes the approximation quantitative: the annealed
score is $C^{\infty}$ in the input, and the gap to the crisp defuzzification is
controlled explicitly by $\beta$ and by how decisively the rules separate.
Together, Theorems~\ref{thm:disjoint}--\ref{thm:anneal} deliver the continuity
and differentiability the action requires, which was the motivating request.

\section{Control: an exact suboptimality certificate}
\label{sec:control}

\subsection{The degenerate case}

\begin{theorem}[TSK reproduces LQR exactly]
\label{thm:c1}
For an LTI plant with quadratic cost, an order-$1$ TSK controller with any
partition of unity and consequents $f_i(x)=-Kx$ for all $i$ satisfies
$u(x)=\sum_i\varphi_i(x)f_i(x)=-Kx$, the LQR optimum, identically.
\end{theorem}

The proof is $\sum_i\varphi_i\equiv1$. We record Theorem~\ref{thm:c1} because
it is the correct baseline and because it is \emph{vacuous as a recommendation}:
it says the model class contains the linear optimum, so any benchmark on which
LQR is optimal cannot discriminate. This is why \S\ref{sec:benchmark} scores on
settling time, peak force and energy jointly---objectives under which LQR is
not optimal---rather than on the quadratic cost it was designed to minimize.

\subsection{The exact certificate, for any convex cost}

Consider a control-affine plant $\dot x=f(x)+g(x)u$ with running cost
$\ell(x,u)=q(x)+c(u)$, optimal value $V^{*}$ and optimal policy $u^{*}$.

\begin{theorem}[Bregman suboptimality identity]
\label{thm:c2prime}
Let $c$ be convex and differentiable, let $V^{*}$ be differentiable and satisfy
the Hamilton--Jacobi--Bellman equation, and let $u(\cdot)$ be admissible (the
closed loop converges to the origin with $V^{*}(x(t))\to0$). Then
\begin{equation}
\boxed{\ J(x_0)-V^{*}(x_0)=\int_0^{\infty}
D_c\big(u(t)\,\|\,u^{*}(x(t))\big)\,dt\ }
\label{eq:c2prime}
\end{equation}
where $D_c(u\|v)=c(u)-c(v)-c'(v)(u-v)$ is the Bregman divergence of $c$.
\end{theorem}

\begin{proof}
HJB stationarity gives $c'(u^{*})=-g^{\top}\nabla V^{*\top}$. Hence
\[
\ell(x,u)+\nabla V^{*}(f+gu)
=q+c(u)+\nabla V^{*}f+\nabla V^{*}gu
=c(u)-c(u^{*})-c'(u^{*})(u-u^{*})=D_c(u\|u^{*}),
\]
using $q+\nabla V^{*}f+\nabla V^{*}gu^{*}+c(u^{*})=0$ (HJB). The left side is
$\ell+\frac{d}{dt}V^{*}(x(t))$ along the closed loop; integrating from $0$ to
$\infty$ and using admissibility to kill the boundary term gives
\eqref{eq:c2prime}.
\end{proof}

\begin{corollary}
For $c(u)=u^{\top}Ru$, $D_c(u\|u^{*})=(u-u^{*})^{\top}R(u-u^{*})$ and
\eqref{eq:c2prime} reduces to the familiar
$J-V^{*}=\int(u-u^{*})^{\top}R(u-u^{*})dt$.
\end{corollary}

Three features deserve emphasis. The identity is \emph{exact}---no constants, no
inequalities. Non-negativity of the gap is free, because $D_c\ge0$ \emph{is}
convexity of $c$. And it holds for costs that are not polynomial: we verified
it numerically for $c(u)=\cosh u-1$ and $c(u)=u^2+u^4$ on inverse-optimal
benchmarks, with residuals of $10^{-11}$ to $10^{-12}$
(Table~\ref{tab:bregman}).

\begin{table}[t]
\centering
\caption{Numerical check of \eqref{eq:c2prime} on inverse-optimal benchmarks.}
\label{tab:bregman}
\begin{tabular}{llccc}
\toprule
cost & controller & gap & $\int D_c\,dt$ & residual \\
\midrule
$\cosh u-1$ & $0.5u^{*}$ & 3.099e$-$1 & 3.099e$-$1 & 3.2e$-$12 \\
$\cosh u-1$ & $1.7u^{*}$ & 5.172e$-$1 & 5.172e$-$1 & 4.0e$-$11 \\
$\cosh u-1$ & $-0.4x$    & 7.536e$-$1 & 7.536e$-$1 & 2.8e$-$11 \\
$u^2+u^4$   & $0.5u^{*}$ & 1.857133e$-$1 & 1.857133e$-$1 & 1.1e$-$11 \\
$u^2+u^4$   & $-0.4x$    & 8.955264e$-$2 & 8.955264e$-$2 & 1.4e$-$11 \\
\bottomrule
\end{tabular}
\end{table}

\paragraph{The hypotheses bind.} Admissibility is not decorative. A controller
fitted without the constraint $u(0)=0$ leaves a steady-state offset, the closed
loop settles to a nearby point rather than the origin, $V^{*}(x(t))\not\to0$,
and the identity silently fails. We encountered exactly this, and it is also the
reason the constraint appears in \S\ref{sec:sos}.

\paragraph{What lies outside.} Theorem~\ref{thm:c2prime} certifies integral
functionals of the closed loop. Peak force $\max_t|u|$ and settling time are not
integral functionals---one is a supremum, the other an exit time---so two of the
three benchmark objectives in \S\ref{sec:benchmark} are provably outside the
certified envelope, and are reported as measurements.

\section{Benchmark: two-mass-spring}
\label{sec:benchmark}

\subsection{Setup}

Two unit masses coupled by a unit spring, force applied to the first mass,
saturating at $|u|\le1$; state $z=(x_1,x_2,v_1,v_2)$. Controllers are scored by
the equal-weight normalized scalarization
\begin{equation}
\mathrm{score}=\tfrac13\left(
\frac{T_{\text{settle}}}{T^{\text{ref}}}
+\frac{\max_t|u|}{|u|^{\text{ref}}}
+\frac{\int u^2dt}{E^{\text{ref}}}\right),
\end{equation}
with the reference the best LQR over a gain sweep, so LQR scores $1.000$ by
construction. The question asked was how fast the least-action fit converges
\emph{with no refinement}: the recipe is applied once, with no tuning---weighted
$k$-means rule placement on open-loop optimal trajectories, then variable
projection of the consequents (one linear solve, globally optimal by
Theorem~\ref{thm:varpro}), then simulate.

\subsection{It converges in two rules}

\begin{table}[t]
\centering
\caption{Convergence with rule count. ``shift'' is the maximum distance from a
closed-loop state to the nearest training sample, in units of the training
set's per-axis spread.}
\label{tab:convergence}
\begin{tabular}{lrrrrrrc}
\toprule
rules & params & settle & peak $|u|$ & energy & score & shift & settled \\
\midrule
LQR            & 4  & 19.797 & 1.0000 & 0.5375 & 1.0000 & --- & yes \\
open-loop ref  & -- & 10.972 & 0.8615 & 0.5585 & 0.8183 & --- & yes \\
1  & 5  & 26.773 & 0.9575 & 0.4214 & 1.0313 & 3.05 & yes \\
\textbf{2}  & \textbf{10} & 17.830 & 0.9088 & 0.4248 & \textbf{0.8666} & 3.07 & yes \\
\textbf{3}  & \textbf{15} & 17.217 & 0.9023 & 0.4390 & \textbf{0.8629} & 3.05 & yes \\
4  & 20 & 18.470 & 0.9088 & 0.4465 & 0.8908 & 3.07 & yes \\
6  & 30 & 18.257 & 0.8560 & 0.4798 & 0.8903 & 3.07 & yes \\
8  & 40 & 19.040 & 0.5531 & 0.6121 & $\infty$ & 3.07 & \textbf{no} \\
12 & 60 & 14.313 & 1.0000 & 1.2041 & 1.3211 & 4.67 & yes \\
16 & 80 & 34.990 & 1.0000 & 7.8944 & $\infty$ & 6.73 & \textbf{no} \\
\bottomrule
\end{tabular}
\end{table}

Ten parameters take the score from $1.031$ (worse than LQR) to $0.867$; the
third rule adds $0.4\%$ and that is the end of it. Two rules close $75\%$ of the
gap between the best LQR and the open-loop reference,
$(1.000-0.863)/(1.000-0.818)$. The win is specific rather than uniform: at
$N=3$ the fuzzy controller uses $18\%$ less energy than the best LQR and $21\%$
less than the open-loop reference it was trained on, and $10\%$ less peak force,
while being \emph{worse} on settling time ($17.2$ against $11.0$). It trades
response speed for actuator economy.

\subsection{Then it degrades, for a diagnosable reason}

Past $N\approx3$ the score degrades monotonically and by $N=8$ the closed loop
stops settling. This is not a numerical artifact and not a failure of the
variable projection, which remains a globally optimal linear solve at every rule
count. It is \textbf{distribution shift}, and the shift column of
Table~\ref{tab:convergence} tracks the failure exactly: $3.05$ while the method
works, $4.67$ at $N=12$, $6.73$ at $N=16$. The controller is fitted only on
optimal trajectories, so it has no data off them; more rules mean each rule is
supported by fewer samples and extrapolates harder. Hence:

\begin{quote}
\textbf{Rule count and closed-loop quality are not monotonically related.} The
least-action fit is globally optimal \emph{for the data it was given}, and the
data is the binding constraint---not the model class and not the optimizer.
\end{quote}

Off-trajectory augmentation (perturbed tubes, or one DAgger round) removes the
collapse but does not move the plateau near $0.86$; nor do higher-order
consequents, which shift it by $1.6\%$ while fit accuracy and closed-loop score
become \emph{inversely} related. Both experiments point at the objective rather
than the model.

\subsection{Direct policy optimization breaks the plateau}

Testing the implied fix---stop fitting sampled targets, optimize the closed-loop
objective directly over the FIS parameters---requires surrendering
Theorem~\ref{thm:varpro}: the closed-loop objective is not quadratic in
$\theta$, so only derivative-free search remains. We use Powell on a shaped
surrogate in which each discontinuous term is replaced by a finite analogue
(soft time outside a tolerance ball; log-sum-exp peak; unchanged energy; plus a
terminal-norm term that makes a non-settling controller \emph{improvable}
rather than merely rejected).

\begin{table}[t]
\centering
\caption{Direct policy optimization, warm-started from the least-action fit,
$600$ closed-loop evaluations each.}
\label{tab:policyopt}
\begin{tabular}{lrrrr}
\toprule
rules & params & imitation & direct & improvement \\
\midrule
\textbf{2} & \textbf{10} & 0.8666 & \textbf{0.6547} & \textbf{24.5\%} \\
3 & 15 & 0.8629 & 0.6729 & 22.0\% \\
4 & 20 & 0.8908 & 0.7565 & 15.1\% \\
\bottomrule
\end{tabular}
\end{table}

At two rules the optimized controller improves \emph{all three objectives at
once}---$31\%$ faster settling, $35\%$ less peak force and $37\%$ less energy
than the best LQR---which no imitation controller did. It is $34.5\%$ better
than the best LQR, $22.9\%$ better than the best imitation result, and $20\%$
better than the open-loop reference.

\paragraph{The warm start is load-bearing.} The same structure started from zero
consequents instead of the least-action fit failed outright: $601$ evaluations,
no stabilizing controller found. The two stages are therefore complementary
rather than competing.

\begin{quote}
\textbf{The least-action fit is not the final method---it is the initializer
that makes the final method reachable.} On its own it caps out at $0.863$;
without it, direct optimization does not start.
\end{quote}

This rehabilitates variable projection in a specific role: its value is not that
it produces the best controller, but that it produces a \emph{stabilizing} one
in a single globally optimal linear solve, with no search and no initialization
problem of its own.

\section{Stability certification by sum of squares}
\label{sec:sos}

\subsection{Rational membership functions}

Replacing Gaussians with $\pi$-shaped (Cauchy) membership functions
$\mu_i(z)=\big(1+\|(z-a_i)/b_i\|^{2}\big)^{-1}$ makes every $\mu_i$ rational,
hence $\varphi_i$ rational, hence the TSK control law exactly a ratio of
polynomials $u(z)=P(z)/Q(z)$ with $Q>0$ structurally. The closed loop is then a
rational vector field and the Lyapunov decrease condition becomes a polynomial
inequality after clearing denominators---which is what a sum-of-squares (SOS)
program can decide.

\subsection{The certificate}

With $V(z)=z^{\top}Pz$, we certify $\Vdot<0$ on the sublevel set
$\{V\le\rho\}\setminus\{0\}$ via the S-procedure: find $\sigma$ SOS with
\begin{equation}
W(z)-\varepsilon\|z\|^{2}-\sigma(z)\big(\rho-V(z)\big)\ \text{is SOS},
\label{eq:sproc}
\end{equation}
where $W=-\Vdot\cdot Q^{2}$ is the numerator-cleared decrease condition. Two
implementation facts were decisive.

\paragraph{The constraint $u(0)=0$ is not cosmetic.} An unconstrained fit gives
$u(0)=O(10^{-4})$, so the origin is \emph{not} an equilibrium, the closed loop
settles to a nearby offset point, and no Lyapunov function can certify
asymptotic stability to a non-equilibrium. The same defect voids
Theorem~\ref{thm:c2prime} through its admissibility hypothesis. Because $u(0)=0$
is a single linear equality in $\theta$, it is imposed exactly by a KKT system,
so the fit remains one linear solve and remains the \emph{global} optimum
subject to the constraint---Theorem~\ref{thm:varpro} survives intact.

\paragraph{Both SOS bases must start at degree $1$.} At $z=0$ the target of
\eqref{eq:sproc} is $-\sigma(0)\rho<0$, which cannot be SOS unless
$\sigma(0)=0$. Left unconstrained, the solver silently drives $\sigma(0)\to0$
and every Gram matrix sits on the positive-semidefinite boundary at $10^{-10}$.
Forcing both bases to begin at degree $1$ gives a genuine margin of
$3.5\times10^{-3}$ at the same certified $\rho$.

\subsection{The obstruction was the solver}

\begin{table}[t]
\centering
\caption{Gram minimum eigenvalue by solver and multiplier degree
($\checkmark$ = certified).}
\label{tab:solver}
\begin{tabular}{llccc}
\toprule
solver & $\sigma$ degree & $\rho=0.5$ & $\rho=0.8$ & $\rho=1.0$ \\
\midrule
SCS (first-order)        & 4 & $-0.158$ & $-0.236$ & $-0.250$ \\
\textbf{CLARABEL} (interior-point) & \textbf{2} & $-1.7$e$-10$ $\checkmark$ & $-1.6$e$-10$ $\checkmark$ & $-4.9$e$-11$ $\checkmark$ \\
\textbf{CLARABEL} (interior-point) & \textbf{4} & $-8.2$e$-10$ $\checkmark$ & $-1.1$e$-09$ $\checkmark$ & $-1.8$e$-10$ $\checkmark$ \\
\bottomrule
\end{tabular}
\end{table}

CLARABEL closes at multiplier degree $2$---the degree at which SCS reported a
violation of $-0.25$---in well under a second. The diagnostic that identified
this was the \emph{relative} violation $|\lambda_{\min}|/\max|G|$: SCS bottoms
out at $10^{-6}$ to $10^{-5}$, a first-order method's realistic accuracy on this
problem, while CLARABEL reaches $10^{-11}$. Reading the absolute eigenvalue
suggested gross infeasibility; the relative figure showed it was noise. We
record this because it cost four hours and because the conclusion generalizes:
on small, badly scaled SOS programs, interior-point and first-order solvers are
not interchangeable.

The certified region for the two-rule imitation controller is
$\{z^{\top}Pz\le1.1144\}$.

\subsection{Exact rational verification}

An interior-point SDP certifies only to solver tolerance. We upgrade the result
by the standard rounding-and-projection procedure (Peyrl and Parrilo): solve
with a strictly positive margin; round both Gram matrices to rationals with
bounded denominator; rebuild the target polynomial exactly over $\mathbb{Q}$
(IEEE doubles \emph{are} rationals, so only the deliberate rounding introduces
error); project the rounded Gram onto the affine set where the polynomial
identity holds exactly, which decouples monomial-by-monomial and is available in
closed form; and verify positive definiteness in exact arithmetic.

Two practical notes. Rational $LDL^{\top}$ is exact but useless at this size,
because denominators square at each elimination step and the $14\times14$
factorization does not finish; integer leading principal minors computed by
fraction-free (Bareiss) elimination have no such blow-up. And
\texttt{nsimplify(..., rational=True)} must be avoided---it searches for
``nice'' closed forms and returns objects such as $2^{818/971}3^{651/971}$ for
an ordinary double---in favour of the exact dyadic conversion.

The result: at $\rho=1/2$ the projected identity has residual exactly $0$, and
all $14$ leading principal minors of the $14\times14$ Gram are positive
integers, the largest a $27{,}737$-bit integer. Therefore
\emph{$z^{\top}Pz$ decreases strictly on $\{z^{\top}Pz\le1/2\}\setminus\{0\}$
for the rationalized controller}, proved rather than measured, with no
floating-point step remaining in the chain.

\subsection{Certification does not scale in rule count}

Since $\deg Q=2(N-1)$, the SDP grows quickly: the Gram is $35\times35$ at
$N=2$, $70\times70$ at $N=3$, $126\times126$ at $N=4$ and $715\times715$ at
$N=8$. $N=2$ solves in seconds; $N=8$ is out of reach with these solvers.
Certification is tractable only at low rule count. It happens to land well here,
since \S\ref{sec:benchmark} found convergence at two or three rules, but any
future move toward more rules trades certifiability for performance.

\section{Certificate-aware policy optimization}
\label{sec:aware}

\subsection{An uncontrolled trade}

Certifying the directly optimized controller of \S\ref{sec:benchmark} exposed a
trade that nothing in the objective knew about: performance improved
($1.044\to0.773$) while the certified ball \emph{shrank} ($0.568\to0.274$). The
optimizer buys performance by acting more aggressively, and aggression is
exactly what a Lyapunov certificate charges for.

\subsection{A proxy that is safe to optimize against}

Putting the certified radius into the objective is obstructed by cost: one SOS
solve is ${\sim}1.5$\,s and a bisection over two Lyapunov candidates is $47$\,s
per controller, against a $601$-evaluation search. We therefore optimize a
sampled surrogate,
\begin{equation}
\hat\rho=\min\{\,z^{\top}Pz\ :\ \Vdot(z)\ge0\,\},
\qquad \hat r=\sqrt{\hat\rho/\lambda_{\min}(P)},
\label{eq:proxy}
\end{equation}
estimated over a fixed sample set. Its logical status is what makes this sound:
a state with $\Vdot(z)\ge0$ \emph{proves} the SOS program must fail at that
$\rho$, since a certificate is a statement about every point of the sublevel
set. Hence $\hat\rho$ is an \emph{upper bound} on any achievable certified
$\rho$---a necessary condition, never sufficient. That is the right direction
for a search objective and the wrong direction for a claim, so the proxy is
optimized against and never reported: every controller is subsequently passed
through the real semidefinite program.

The penalized objective is
\begin{equation}
J_w(\theta)=\mathrm{cost}(\theta)
+w\,\frac{\max\!\big(0,\,r_t-\hat r(\theta)\big)^{2}}{r_t^{2}},
\label{eq:penalty}
\end{equation}
a one-sided squared hinge---exceeding the target is never punished---normalized
so that $w$ is dimensionless and comparable to the score. The search remains
restricted to the $u(0)=0$ null space of \S\ref{sec:sos}.

\subsection{The front, and a result in the opposite direction}

\begin{table}[t]
\centering
\caption{Certificate-aware policy optimization. ``SOS ball'' is the real
certificate, best of two Lyapunov candidates; ``held out'' is the score on six
initial conditions not used in the objective.}
\label{tab:front}
\begin{tabular}{lrrrrl}
\toprule
$w$ & score & held out & proxy $\hat r$ & SOS ball & $V$ from \\
\midrule
--- (imitation) & 1.0437 & 1.1395 & 0.6718 & \textbf{0.5684} & Riccati \\
0               & 0.7725 & 0.8272 & 0.3900 & 0.2729 & lin \\
0.1             & 0.8155 & 0.7717 & 0.8676 & 0.3452 & lin \\
\textbf{0.3}    & \textbf{0.7713} & \textbf{0.7194} & 0.6124 & \textbf{0.3354} & lin \\
1               & 0.7906 & 0.7230 & 0.7039 & 0.3498 & lin \\
3               & 0.8265 & 0.8105 & 0.8676 & \textbf{0.3579} & lin \\
\bottomrule
\end{tabular}
\end{table}

The expected outcome was a strict trade. Instead, $w=0.3$ \textbf{dominates
$w=0$ on all three axes}: better score ($0.7713$ against $0.7725$), better
held-out score ($0.7194$ against $0.8272$), and a larger certified ball
($0.3354$ against $0.2729$). The held-out column explains why. The unpenalized
search overfits its six training initial conditions, degrading $7\%$ out of
sample, and the certificate penalty suppresses precisely the aggressive
behaviour responsible.

\begin{quote}
\textbf{A Lyapunov certificate acts as a regularizer.} It charges for the same
aggression that costs generalization, so the first increment of
certificate-awareness is free.
\end{quote}

Past that point the trade is real and $w$ prices it: over $w\in\{0.3,1,3\}$ the
certified ball rises $0.3354\to0.3498\to0.3579$ while the score degrades
$0.7713\to0.7906\to0.8265$.

\section{Calibrating the proxy}
\label{sec:calibrate}

\S\ref{sec:aware} left the certificate a dial rather than a specification: the
proxy over-stated the certified ball by $1.18$ to $2.51\times$, so asking for
$0.67$ delivered $0.35$. We had attributed this to S-procedure conservatism,
which would have required a higher multiplier degree. It was instead
\emph{estimator error}, of two distinct kinds.

\subsection{Two defects, both in the estimator}

\paragraph{A minimum-order statistic.} Violations are not rare---$37\%$ of
cloud points satisfy $\Vdot\ge0$---so the sampled minimum in \eqref{eq:proxy}
depends entirely on whether some point happened to land near the inner
boundary. Across five seeds it reports $0.6718$, $0.6744$, $0.7484$, $0.6608$,
$0.7369$: a $12.5\%$ spread. Fixing the seed made it deterministic, not
accurate.

\paragraph{The obvious repair is degenerate.} Sharpening \eqref{eq:proxy} by
locally solving $\min\{z^{\top}Pz:\Vdot(z)\ge0\}$ fails structurally rather than
numerically: \textbf{$z=0$ is feasible}, since $\Vdot(0)=0$ exactly. The program
has global optimum $0$ at the origin and a descent method walks to it. Measured:
of six local solves from the best sampled starts, two collapsed to $\|z\|=0$ and
all six terminated infeasible on the $\Vdot=0$ boundary, so a strict feasibility
recheck discards every one and the refined bound never moves.

\subsection{Ray search}

Both defects vanish if the search is organized by direction. Along a fixed unit
vector $d$ define the first violation radius
\begin{equation}
r(d)=\min\{\,r>0:\Vdot(rd)\ge0\,\},\qquad
\hat\rho=\min_{d}\;r(d)^{2}\,d^{\top}Pd,
\label{eq:ray}
\end{equation}
computed by scan-and-bisect over a fixed direction set. There is no optimizer,
no degenerate solution, and the origin is excluded by construction because the
scan starts at $r>0$. Bisection returns the \emph{outer} bracket endpoint, which
is a state genuinely satisfying $\Vdot\ge0$, so the necessary-condition property
that makes the proxy safe to optimize against is preserved exactly. This is
checked rather than argued: all $973$ returned endpoints satisfy $\Vdot\ge0$ on
re-evaluation, and the binding one is the explicit state
$z^{*}=(+0.0237,+0.0778,+0.0628,-0.2563)$ with $\Vdot=+2.9\times10^{-9}$.

\begin{table}[t]
\centering
\caption{Estimator comparison, same controller and same necessary condition.}
\label{tab:estimator}
\begin{tabular}{lccccccc}
\toprule
estimator & seed 0 & seed 1 & seed 2 & seed 3 & seed 4 & spread & vs SOS \\
\midrule
cloud & 0.6718 & 0.6744 & 0.7484 & 0.6608 & 0.7369 & 12.5\% & $1.182\times$ \\
\textbf{ray} & \textbf{0.5937} & \textbf{0.5934} & \textbf{0.5849} &
\textbf{0.5937} & \textbf{0.5915} & \textbf{1.5\%} & $\mathbf{1.044\times}$ \\
\bottomrule
\end{tabular}
\end{table}

\subsection{What remains is a constant}

Perturbing the fit along the $u(0)=0$ null space produces controllers spanning a
range of aggressiveness, each still with an equilibrium at the origin. Over
five that certify, the ratio of the true certificate to the ray estimate is
\begin{equation}
\kappa=\operatorname{median}\!\left(\frac{\text{SOS}}{\text{ray}}\right)=0.9620,
\qquad\text{range }[0.9271,\,0.9735].
\end{equation}
Of the $3.8\%$ shortfall, $2.5$ points are the deliberate $0.95$ shrink applied
to $\rho$ before re-certifying ($\sqrt{0.95}=0.9747$). Genuine degree-$4$
S-procedure conservatism is therefore about $1.3\%$---confirming that chasing
multiplier degree would have bought nothing.

\subsection{A specified radius, confirmed by the solver}

Hinging \eqref{eq:penalty} on the calibrated estimate $\kappa\hat r$ states the
target directly in certified-ball units. A pass that misses raises the target by
exactly the shortfall; a pass that misses \emph{and} fails to improve on the
previous one is not short of target but stuck, and the target is pushed $50\%$
further.

\begin{table}[t]
\centering
\caption{Specifying a certified radius. Every ``yes'' is a semidefinite-program
certificate, not a proxy estimate.}
\label{tab:spec}
\begin{tabular}{ccrrrrrrc}
\toprule
spec & pass & score & held out & $\kappa\hat r$ & SOS Ric & SOS lin & best & met? \\
\midrule
0.35 & 1 & 0.7956 & 0.7546 & 0.3478 & 0.0000 & 0.3415 & 0.3415 & no \\
0.35 & 2 & 0.7962 & 0.7227 & 0.3577 & 0.0000 & 0.3549 & \textbf{0.3549} & \textbf{yes} \\
0.45 & 1 & 0.7690 & 0.7296 & 0.5277 & 0.0000 & 0.3452 & 0.3452 & no \\
0.45 & 2 & 0.7857 & 0.7536 & 0.5865 & 0.0000 & 0.3430 & 0.3430 & no \\
0.45 & 3 & 0.8487 & 0.8667 & 0.7269 & 0.6782 & 0.3654 & \textbf{0.6782} & \textbf{yes} \\
0.55 & 1 & 0.7717 & 0.7466 & 0.5503 & 0.1038 & 0.3368 & 0.3368 & no \\
0.55 & 2 & 0.8469 & 0.8560 & 0.7278 & 0.6791 & 0.0000 & \textbf{0.6791} & \textbf{yes} \\
\bottomrule
\end{tabular}
\end{table}

All three specifications are met. The best controller of the entire study falls
out of this procedure: \textbf{certified ball $0.6791$ at score $0.8469$}, which
dominates the imitation controller on \emph{both} axes---a larger certified
region than $0.5684$ and a $19\%$ better score than $1.0437$---and nearly
doubles the best certified region obtained in \S\ref{sec:aware}.

\subsection{Two limits calibration does not remove}

\paragraph{A single $\kappa$ can fail.} At specification $0.45$, pass $1$,
$\kappa\hat r=0.5277$ against a certificate of $0.3452$: a ratio of $0.65$, far
outside $[0.927,0.974]$. The per-candidate columns show why---$\text{SOS
Ric}=0.0000$. The proxy takes a maximum over Lyapunov candidates and attains it
at the Riccati matrix, where $\Vdot<0$ genuinely does hold that far out, but
which the degree-$4$ S-procedure cannot certify \emph{at all}. The proxy
measures where $\Vdot<0$ \emph{is true}; SOS measures where it is \emph{provable
at this degree}; the argmax of the two differs. Calibration handles the gap when
both agree on the candidate, and the outer loop absorbs it when they do not.

\paragraph{The achievable set is bimodal.} Every row lands near either $0.35$ or
$0.68$, according to which Lyapunov function certifies---two basins, not a
continuum. Specification $0.45$ therefore fails twice and then succeeds, and
asking for $0.45$ delivers $0.6782$: the specification is \emph{met, not
matched}. This is a property of the landscape, not of the calibration. Both
limits point at the same remedy: searching $V$ jointly with $\theta$, a bilinear
problem, would remove the max-over-two-candidates artifact that produces the
calibration outlier and the bimodality alike.

\section{What is proved, and what is not}
\label{sec:inventory}

Table~\ref{tab:inventory} separates the results by evidentiary status:
analytically proven, certified numerically, or unavailable. We regard keeping
this distinction visible as part of the method rather than as bookkeeping,
because the value of a variational formulation lies precisely in which claims it
can discharge without a solver.

\begin{table}[t]
\centering
\caption{Guarantee inventory.}
\label{tab:inventory}
\footnotesize
\begin{tabular}{@{}ll>{\raggedright\arraybackslash}p{3.1cm}>{\raggedright\arraybackslash}p{5.2cm}@{}}
\toprule
\# & Guarantee & Status & Conditions \\
\midrule
G1 & Consequents globally optimal, closed form & Proven & antecedents fixed; independent regressors \\
G2 & Action convex in function space & Proven & $\lambda\ge0$ \\
G3 & Reduced gradient exact & Proven & $\delta$-coverage \\
G4 & Greedy $=$ joint iff Gram block-diagonal & Proven & for all targets \\
G5 & $\lambda^{*}=3b^{4}/(8c^{2})$ & Proven & 2 equal-width Gaussians \\
G6 & $\vee=+$ iff disjoint output supports & Proven & --- \\
G7 & Centroid $=$ order-$0$ TSK & Proven & disjoint, product implication \\
G8 & Annealing error bound & Proven & --- \\
G9 & TSK reproduces LQR exactly & Proven & LTI, quadratic cost (and vacuous) \\
G10 & Bregman suboptimality identity & Proven & control-affine, convex $c$, admissible, $V^{*}$ known \\
G11 & Local optimality of antecedents & Certified numerically & KKT on critical cone $+$ FD Hessian \\
G12 & Stability / region of attraction & Proven by SOS & rational MFs, $u(0)=0$, interior-point SDP \\
G13 & Global optimality over antecedents & Not available & --- \\
G14 & Certificate for settling time, peak force & Provably outside G10 & not integral functionals \\
\bottomrule
\end{tabular}
\end{table}

Two obstructions did not yield and should be stated plainly. G13 is genuinely
open: the antecedent landscape is multimodal and nothing here bounds the gap to
the global optimum. G14 is not an open problem but a structural exclusion---a
supremum and an exit time are not integral functionals, so no argument of the
form of Theorem~\ref{thm:c2prime} can reach them. The honest route to certifying
the benchmark objectives is to \emph{restate} them as integral costs
($\int\phi(u)dt$ with $\phi$ a convex barrier; $\int w(t)\|x\|^{2}dt$), which
changes the problem in a direction the framework can certify.

\section{Machine verification}
\label{sec:verification}

Three tiers of checking were applied, weakest to strongest.

\textbf{Tier 1: symbolic identity checking.} Twenty-one identities---the
Euler--Lagrange equation, Galerkin orthogonality, the envelope-theorem gradient,
the $\lambda^{*}$ integrals, the Bregman identity, the partition-of-unity
consequences---are verified as symbolic identities in a computer algebra system
rather than at sampled points. All $21$ pass. (One subtlety: a matrix comparison
\texttt{Matrix(...) == 0} is \texttt{False} even for the zero matrix; using the
correct predicate moved the count from $19/21$ to $21/21$.)

\textbf{Tier 2: exact rational proof of the SOS certificate}, as described in
\S\ref{sec:sos}: residual exactly zero, all $14$ integer leading principal
minors positive.

\textbf{Tier 3: proof assistant.} A Lean/Mathlib formalization would close the
gap Tier 1 leaves---the analytic hypotheses---but could not be attempted here,
as the toolchain host is unreachable under this environment's network policy.
Rather than commit unverifiable source we record where it would pay.
Theorem~\ref{thm:c2prime} is the clear candidate: its content is an
integration-by-parts argument plus a limit, and its hypotheses (admissibility,
convexity, differentiability of $V^{*}$) are exactly the kind that get quietly
dropped---as \S\ref{sec:control} notes, we dropped one empirically. The SOS side
should stay in exact integer arithmetic, where it already is: formalizing it
would mean formalizing Bareiss determinants to re-derive an answer that is
already decidable and decided.

\section{Limitations and future work}

\begin{itemize}
\item \textbf{One plant, one dimension of universe.} The approximation results
are scalar-input; $\lambda^{*}$ in the multi-input case turns the crossover into
a hypersurface with a direction-dependent width and may have no clean form. The
control results use a single benchmark.

\item \textbf{Certification does not scale.} At $N=8$ the SDP is $715\times715$
and out of reach with these solvers. This bounds the whole certified story to
low rule count.

\item \textbf{Two fixed Lyapunov candidates.} Every certified radius reported is
a maximum over the Riccati matrix and the closed-loop linearization. Searching
$V$ jointly with $\theta$ is bilinear, was not attempted, and is the single
change that would most improve \S\ref{sec:calibrate}.

\item \textbf{The classifier results have not met real data.}
Theorems~\ref{thm:disjoint}--\ref{thm:anneal} are proved but exercised only on
synthetic partitions. This does not test the guarantees, which are already
proven; it tests whether they bind on anything real.

\item \textbf{Relative form of Theorem~\ref{thm:c2prime}.} The certificate
presupposes $V^{*}$, which is what a forward problem lacks. For any $\hat V$
satisfying HJB with residual $\varepsilon$, the same algebra should bound the
true gap in terms of $\varepsilon$; the useful part of that derivation is
determining in which norm $\varepsilon$ must be measured. Without it, the
certificate is restricted to inverse-optimal benchmarks---problems whose answer
was known in advance.

\item \textbf{Analytic Hessian.} G11 is certified numerically only because the
Hessian is obtained by finite differences. The same envelope argument
differentiated once more would upgrade it.
\end{itemize}

\section*{Reproducibility}

All results are produced by a deterministic, seeded implementation; recorded
output is stored alongside the source, and the numerical constants that matter
(quadrature order, ridge scaling, optimizer schedule, finite-difference steps,
definiteness tolerances) are documented individually together with the failure
each was chosen to prevent. Tables~\ref{tab:convergence}, \ref{tab:policyopt},
\ref{tab:front}, \ref{tab:estimator} and \ref{tab:spec} are transcribed directly
from that recorded output.

\end{document}
